% latexmk -cd -lualatex works/L_1/appendix.tex
% latexmk -c -cd works/L_1/appendix.tex

\documentclass{ees}

\pghyphenation{english}{Er-götz-lich-keit Tändel-marckht}

\newlist{lyricslist}{description}{1}
\setlist[lyricslist]{
  partopsep=0pt,
  labelindent=0em,
  labelwidth=4.5em,
  leftmargin=4.5em,
  labelsep=0pt,
  first=\ltseries,
  font=\normalfont\itshape\ltseries,
  style=multiline
}

\newenvironment{lyrics}[1]{%
  \subsection{#1}\nopagebreak%
  \begin{lyricslist}%
  \let\voice\item%
}{%
  \end{lyricslist}%
}

\begin{document}

\clearscrheadfoot

\chapter{Prefatory pages}

\textit{[bc part, fol. 1r]}\\
Zwey neue | und extra luſtige muſicaliſche | Tafel-Stücke. | I. | Der Wienneriſche Tändlmarckt. | Mit 4. Singſtimmen, 2. Violinen und Baſſo ordinario. | Die Singſtim- | men ſtellen vor: Den Sägfeiller, Hollhippen-Kramer, einen Savoyarden | mit ſeinem Spielwerck, und einen Marcktſchreyer. Die Inſtru- | menten aber imitiren das ſogenannte Winne- | riſche Breinglöckl. | II. | Die Bauren-Richters-Wahl. | Mit 5. Singſtimmen. 2. Violinen und Baſſo ordinario. Die Singer ſind: | Der Bartlme Zimmermann, als Senior, der Paul Schnepfendreck, der | Böhmiſche Mazko, der Schwäbiſche Wirth, und der einäugige | Schulmeiſter, welche ein jeder in ſeiner Sprache, | die Vota zur Wahl geben. | Beyde Stücke/ ſind nach Art des Quodlibete eingerich- | tet, und jedes, zur angenehmen Abwechslung, mit 7. Arien, und darzu | gehörigen Recitativen verſeheh, mithin bey vornehmen Gaſt- | mahlen und anderen luſtigen Geſellſchaften mit | groſſem Applauſu zu produciren. | Von dem Autore des Muſicaliſchen Calenders | Gregorio Joſepho Werner, | Hochfürſtlich Eſtorhaſiſchen Capellmeiſtern zu Eiſenstadt in Ungarn. | AUGSPURG, | Verlegt von Johann Jacob Lotters ſeel. Erben.

\textit{[fol. 2r]}\\
Avertiſſement. | Hier werden einer courieuſen Welt/ zur Luſt und | Zeit-vertreiblichen Ergötzlichkeit zwey Muſicali- | ſche Quodlibetica vorgeſtellet. Daß erſtere | betitult ſich/ der Wieneriſche Tändelmarckht/ | allwo in dem Choro alles natürlich begriffen/ | wie jedwe- | dere ihre Waaren allda zum Verkauf pflegen auszuruffen/ | worunter das ſogenannte Breinglöckel in denen Inſtru- | mental-Stimmen Imitations-Weiß ſich hören läßt. Als- | dann erfolget Recitativ- und Arien-mäßig der Sägfeil- | ler/ Savojard, Hollhippen-Kramer/ ſamt dem Arzt | auf der Freyung. Alles iſt leicht zu tractiren/ | anbey mit

\textit{[fol. 2v]}\\
guten Fleiß kurz verfaſſet, damt das Auditorium immer | eine angenehme Veränderung hören möge. Das zweyte | Stuck/ wird nach allhießiger Landsart benamßet die Bon- | thätung/ oder die Bauren-Richters-Wahl/ ſo beſtehet | aus 5. Perſonagen; als der Barthlme Zimmermann/ | der Paul Schnepfendreck/ der Schwabnwürth/ der | Böhmiſch Mazko, und der Schulmeiſter/ welcher mit | ſeiner Schul-Jugend zu buchſtabieren anfanget. Alsdann | wird zur Wahl geſchritten/ mit einer vorgehenden Bauren- | Oration, da endlich nach langem Wortſtreit und Knittel- | haften Expreſſionen der Richter einſtimmig erwählet wird. | In allen beyden Pieçen wird ſodann jedem Vocaliſten/ | in ſeiner Perſon hauptſächlich die Pronunciation wohl | anrecommendiret. | Der ich mich übrigens zu der re- | ſpectivè geneigten Liebhaberen Dienſte/ | fernerhin gehor- | ſamſt anerbietend empfehle. | Ein befliſſener Diener | Gregorius Joſeph Werner.

\cleardoublepage
\chapter{Lyrics}

\textit{Roles:} Der Savoiard (T) · Der Hollähippen-Krämer (T) · Arzt (B) · Der Sagfäller (B)
\enlargethispage{2\baselineskip}

\begin{lyrics}{1 Chorus · Das Breinglöckl}
  \voice[Sagfäller]
  Sagfälln

  \voice[Savoiard]
  O ſöner rarite

  \voice[Arzt]
  Sagfälln

  \voice[Savoiard]
  O ſöner Spillwerckh\\
  wer will ſeh guckh in die Lokh.\\
  Käfts Milli Frau än Millräm oder än Budä

  \voice[H.-Krämer]
  Käfts än Sand

  \voice[Sagfäller]
  Bändl käfts än Zwirn

  \voice[Arzt]
  Bändl käfts än Zwirn

  \voice[Sagfäller]
  Sagfälln

  \voice[Savoiard]
  re re re re re re re re re re re re re rettich.

  \voice[Arzt]
  Sagfälln

  \voice[H.-Krämer]
  re re re re re re re re re re re rettich.

  \voice[Sagfäller]
  Kafts Säfäla dri Würſt ſänd gantz warm,

  \voice[Savoiard]
  La laterna magic

  \voice[H.-Krämer]
  La laterna magic

  \voice[Arzt]
  Bändl käfts än Zwirn

  \voice[Sagfäller]
  Bändl käfts än Zwirn

  \voice[Savoiard]
  koäni alti Leibſtücker zu verkäffa\\
  alte Hüet oder Haaſenbalck.

  \voice[alle]
  Aen neue Zeittung um än Creuzer daß Gott ärbarm.
\end{lyrics}

\begin{lyrics}{2 Recitativo}
  \voice[Sagfäller]
  Mei ſchöni Kunſt hat Archimedes erdenckt\\
  und hätt ſi beym Studirn\\
  döß Werckh recht auszuführn\\
  ſchier bald an einer Bradwurſt erhenckt.
\end{lyrics}

\begin{lyrics}{3 Aria}
  \voice[Sagfäller]
  Sagfälln\\
  diſi ſchwäri Kunſt\\
  hab I nit glei umſunſt\\
  in oänen Tag crepirt\\
  biß I erſt glehrnt hab\\
  auf den Tact recht fälln die Sag\\
  in Zwingär äiniſpälln\\
  hat mir den Kopf verwürt\\
  Sagfälln.
\end{lyrics}

\begin{lyrics}{4 Recitativo}
  \voice[Savoiard]
  Ick gom auß die Savoier Länd\\
  mit ſöner Raritè\\
  die Spilliwerg bellè gallant\\
  will ſi vor Creutzer ſeh?\\
  da wird ſi find Bell Jl di Eld\\
  wie ſi mack förkt daß ganſſe Welt\\
  le Cont de Saſs ſis ä ſu ferd\\
  da ſittert unter ſi die Herd,\\
  wer greuſſer ſalt vien preſto qua\\
  Ick ſpill trâtant la Miſuca.
\end{lyrics}

\begin{lyrics}{5 Aria forlana}
  \voice[Savoiard]
  Ô bella Margarita\\
  ô laſcia far pur me\\
  tu Sei il mio Corla vita\\
  jo voglio morir per te\\
  Sa rai la mia Spoſa\\
  ti diro qualche coſa\\
  pe pe perle pe perle pe.
\end{lyrics}

\enlargethispage{\baselineskip}
\begin{lyrics}{6 Recitativo}
  \voice[H.-Krämer]
  Äs iſt än altes Sprichwort\\
  ſo I hab ghert bald da bald dort:\\
  was an den Galgen ghert ertrinckt nit\\
  und was erſauft daß henkt mär nit,\\
  döſtwögn laß I mi gern baden,\\
  weil mir das Waſſer nit kan ſchaden.
\end{lyrics}

\begin{lyrics}{7 Aria}
  \voice[H.-Krämer]
  Allo frisch auf\\
  bräf gwürffelt drauf\\
  I zahl ja rödli auß\\
  änß zwoä löſch ſauber auß\\
  zahl bräf auß\\
  bäm ſchlaperment\\
  s’Krämöl geht drauf\\
  hiezt falt drey vieri drauf\\
  fünft ſechs höbt s’Spill gar auf\\
  zahl braf auß\\
  döß Spill wird ni mär auß\\
  fein ſauber is auß.
\end{lyrics}

\begin{lyrics}{8 Recitativo}
  \voice[Arzt]
  Wohl edle geliebte\\
  verliebt und auch betrübte\\
  wer ſich will helffen ſehn\\
  kan auf die Freyung gehn,\\
  dort hab ich meinen Stand\\
  wie jederman bekandt.\\
  Seyts alle meine Kinder sagts Ja = =
\end{lyrics}

\begin{lyrics}{9 Aria}
  \voice[Arzt]
  Hier hab ich eine Pillen\\
  die wegt biß zehen Pfund,\\
  wer will ſein Kranckheit ſtillen\\
  wird ſchwärlich darvon gſund\\
  ich hab ein Panacea\\
  von Aquafort und Pfeffer,\\
  ſtirbt er daran noch nicht\\
  ſo komt er mir gwiß öffter.
\end{lyrics}

\begin{lyrics}{11 Recitativo}
  \voice[Sagfäller]
  Hiez hab i oänen Vorthl glehrnt\\
  der koſt mi zwey Maß Wein\\
  den ſag i wärlä koänen nit\\
  ſolts ä mei Vadä ſeyn\\
  Dözen krump hin und her zu biegn\\
  iſt gwiß koä Narädey\\
  mär lehrnts nit glei ſo in där Wiegn\\
  äs braucht Värſtand därbey.
\end{lyrics}

\newpage
\begin{lyrics}{12 Aria}
  \voice[Sagfäller]
  Hoch und mitle niedre Stende\\
  brauchen mi und ſänd gar froh,\\
  ſpröchen döß Mann ſeine Hände\\
  ſänd mehr werth als oän Fuhr Stroh,\\
  dößwögn laß I mi bräf zahln\\
  wäll ſö mi vor andern alln\\
  holln und ſagen der verſtehts\\
  alli andri wökher gehts.
\end{lyrics}

\begin{lyrics}{13 Recitativo}
  \voice[Savoiard]
  Allo allegramento!\\
  wer guck in die Laterna magic\\
  da ſeyn ſo vill ſöner Varietè,\\
  o Unvergleiklik luſtig Curioſitè\\
  da gom dö groſſe Mogol\\
  mit undert Elephant\\
  will ſi beſwinge der gänſſe Perſer Länd\\
  da goms di Bader wißi\\
  Mare Antonio Criſtier\\
  ä nide ſo ſtärk a nide ſo ſtärk\\
  ä nide ſo ſtärk ä nide ſo ſtärk\\
  ſonſt dir in die Gſchiſte ofier\\
  da iß ſi meine haltelang gelebte Groß-Vatteri\\
  mit di langi Näß\\
  da la la ra la la ra ra\\
  da la la ra la la ra ra\\
  da ſehe ſi lauteri Späß.
\end{lyrics}

\begin{lyrics}{14 Aria}
  \voice[Savoiard]
  O che bella Tortorella!\\
  O bella Tortorella\\
  che fai ſempre grugeru gru\\
  S’io non amo ti mia donzella\\
  giuro che ſono un becco futu\\
  Sei coſtante al tuo amante\\
  la ſcia paſſare tu mikh ſon verſteh\\
  Sei fedele mamoiſelle\\
  alſo verſpreckh I dir forte di Hee.
\end{lyrics}

\newpage
\begin{lyrics}{15 Recitativo}
  \voice[Arzt]
  Hochgeneigte Compagnie,\\
  ich glaub ihr höret nie\\
  daß mich jemand noch geſchändet\\
  und mein Paquet ſonſt übel hätt\\
  an ſeinen Leib verwendet,\\
  probirt es nur ihr werdet ſehn\\
  daß keiner unter euch allen\\
  mir künfftig wird zu Gfallen\\
  mehr auf die Freyung gehn.
\end{lyrics}

\begin{lyrics}{16 Aria}
  \voice[Arzt]
  Ist etwann ein betagte Frau\\
  mit neunzig Jahr zugegen,\\
  die mach ich jung wann ſie ſchon grau\\
  will ihr die Falten außfegen,\\
  hat einer Wehtag in den Aug\\
  daß er nicht wohl khan ſehen,\\
  der brauch mein ſcharf geſalzne Laug\\
  ſo iſts mit ihm geſchehen.
\end{lyrics}

\begin{lyrics}{17 Recitativo}
  \voice[H.-Krämer]
  Wäll I mei Krämel ganz verſpilt,\\
  ſo iß Hollähippen Brauch\\
  wälß oänen wie den andern gilt\\
  ſo laß I mi taufen auch,\\
  hiez brauchet I halt an Göden,\\
  drum hed I ſchön geböden,\\
  der gnädi Herr wolt willign ein\\
  und hiez mei Bad und Tauf-Gött ſeyn.
\end{lyrics}

\begin{lyrics}{18 Aria}
  \voice[H.-Krämer]
  Äns bitt I mir dö auß\\
  ä Kriegl Wein muß ä da ſeyn\\
  ſonſt kriegät I än Grauß\\
  daß Waſſer ſchmäkt mär ſo nit wohl\\
  där Wein macht mi Guräſchi voll\\
  hiez hoäſts döß Spill iß auß\\
  fey ſauber iß auß.
\end{lyrics}

\end{document}
